% --- Deckblatt Definitionen ---
\def\labtitle{Elektronische Produktentwicklung} % Titel des Labors
\def\labcode{MECH-B-4-AED-EPE1-ILV} % Code des Labors
\def\labname{} % Name des Labors
\def\labnum{Nummer der Laborübung} % Nummer der Laborübung
\def\labdate{Datum Labor} % Datum des Labors
\def\study{Bachelor: Mechatronik Design and Innovation} % Studienrichtung
\def\studygroup{} % Studiengruppe
\def\term{4.} % Semesterzahl
\def\lecturer{Rene Feldkircher} % Dozent
\def\group{BA-MECH-24} % Studiengruppe
\def\student{Seyam Mir Afghan, Elham Mir Afghan, Nils Brieger, Benedikt Motzet} % Studierende
\def\contributor{}
\def\litname{literatur} % Literaturverzeichnis

% --- Kopf- und Fußzeile definieren ---
\fancyhf{} % Alle Kopf- und Fußzeilenfelder leeren
\fancyhead[L]{\labtitle} % Kopfzeile links
\fancyhead[C]{} % Kopfzeile zentriert
\fancyhead[R]{\labname} % Kopfzeile rechts
\renewcommand{\headrulewidth}{0.4pt} % Trennlinie oben
\fancyfoot[L]{\thepage} % Seitenzahl links
\fancyfoot[R]{{\footnotesize\student}} % Studierende rechts
\renewcommand{\footrulewidth}{0.4pt} % Trennlinie unten
\setlength{\headheight}{14.49998pt} % Kopfzeilenhöhe anpassen
\addtolength{\topmargin}{-2.49998pt} % Seitenrand ausgleichen

% --- Tabellenhöhen anpassen ---
\renewcommand{\arraystretch}{1.2} % Zeilenhöhe für Tabellen erhöhen

% --- Bildunterschriften anpassen ---
\captionsetup{belowskip=12pt, aboveskip=4pt} % Abstand über und unter Bildunterschriften
\captionsetup{font=footnotesize} % Schriftgröße für Bildunterschriften

% --- Absatzformatierung ---
\setlength{\parindent}{0em} % Kein Einzug bei Absätzen
\setlength{\parskip}{1.5ex plus0.5ex minus0.5ex} % Abstand zwischen Absätzen

% --- Überfüllungswarnungen minimieren ---
\tolerance = 9999
\sloppy

% --- SI-Einheiten (siunitx Einstellungen) ---
\sisetup{
    output-decimal-marker={,}, % Komma als Dezimaltrennzeichen
    per-mode=reciprocal, % Brüche im Kehrwertmodus
    exponent-product=\cdot, % Malzeichen für Exponenten
    retain-explicit-plus, % Pluszeichen beibehalten
    range-phrase = {\dots}, % Bereich mit "..."
    separate-uncertainty, % Unsicherheit getrennt anzeigen
    list-separator={; }, % Listentrenner
    list-final-separator={; } % Abschließender Listentrenner
}

% --- Diagrammeinstellungen (pgfplots) ---
\pgfplotsset{
    /pgf/number format/use comma, % Komma als Dezimaltrennzeichen
    /pgf/number format/.cd, set thousands separator={\,} % Tausendertrennzeichen
}
